\chapter{绪论}

\section{深度学习求解偏微分方程方法概述}
偏微分方程是计算数学研究的核心问题，在物理学、工程学等多个领域具有广泛应用。传统的数值求解算法，例如有限差分法、有限元方法、有限体积法和谱方法等，往往依赖于基函数或网格划分来进行计算 。近年来，深度神经网络在科学计算领域取得了革命性进展，为偏微分方程的数值解提供了新的思路。2018年以来，随着物理信息神经网络（Physics-Informed Neural Networks, PINNs）和深度里兹法等方法的提出，深度神经网络在求解偏微分方程（Partial Differential Equations, PDEs）方面取得了显著进展。相比于传统的数值方法，如有限元法、有限差分法等，PINNs等新兴方法展示出了独特的优势。

这些新方法的提出，逐渐引起了学术界和工业界的广泛关注。人们开始意识到，深度神经网络能够克服传统数值方法在处理复杂边界条件、多尺度问题以及高维问题时所面临的一些瓶颈和挑战。具体而言，DNNs在逼近复杂函数、处理非线性关系和适应高维空间上表现得尤为出色，这为传统科学计算领域带来了新的思路和工具。


\section{研究现状与相关工作}\label{sec-compile}

目前深度学习求解PDE已经形成了一套较为完善的方法体系和理论分析工具：在算法上，首先是用神经网络代替传统的基函数或网格划分来表示PDE的解，其次是通过设计合适的损失函数来引导神经网络的训练过程，使其满足PDE的物理约束和边界条件，最后是利用梯度下降等优化算法来调整神经网络的参数，从而逼近PDE的解。在理论分析方面，研究者们主要关注神经网络的逼近能力、收敛性以及泛化误差等问题。通过引入Rademacher复杂度、VC维等统计学习理论工具，学者们对深度神经网络在求解PDE中的表现进行了深入分析。此外，一些研究还探讨了不同网络结构（如卷积神经网络、循环神经网络等）在处理特定类型PDE时的优势和局限性。

当前，深度学习在经典线性正问题求解领域已取得显著进展，以Deep Ritz方法为例，其在偏微分方程（PDE）正问题求解框架及相应的误差分析理论方面已趋于成熟。然而，针对特征值问题、非线性椭圆方程及反问题等复杂场景的研究尚显不足。受限于上述问题固有的正交约束、强非线性及不适定性等本质困难，现有框架难以直接迁移应用。因此，构建能够克服上述挑战的新型深度学习求解策略并完善其理论体系，对于拓展该领域的研究边界具有至关重要的意义。

\begin{figure}
\begin{center}
\includegraphics[width=0.8\textwidth]{figures/PDE.png}
\end{center}
\end{figure}
\section{本文的主要工作和贡献}
本文提出了了若干类椭圆型PDE问题的深度学习求解方法。

第二章中，针对高维椭圆算子的多重特征值求解问题，本文提出了一种改进的深度里兹法。本文设计了一种新的包含变分形式、边界惩罚和正交惩罚的损失函数，避免了传统瑞利商方法中的非凸性，证明了该损失函数的极小值点确为原方程的特征函数，给出了详细的误差分解，包括逼近误差、统计误差（利用 Rademacher 复杂度）和累积误差，此外还在10维立方体和10维、20维球体等高维区域上验证了算法的有效性。

第三章中，针对带有 Dirichlet 和 Robin 边界条件的半线性椭圆方程（$-\Delta u + f(u) = g$），本文构建了基于变分原理的深度学习求解框架。利用带有 $\mathrm{ReLU}^2$ 激活函数的残差神经网络（ResNet），构建了基于能量泛函的损失函数，并证明了其极小值点的存在性与唯一性。本文还推导了数值解与真实解之间 $H^1$ 范数误差的显式上界，将总误差分解为逼近误差、统计误差和截断误差，并给出了每一项的收敛阶。数值实验表明，该方法能够有效处理高维空间（如 10 维）中的非齐次边界问题及非光滑解问题。

第四章针对从噪声观测数据中重构椭圆方程未知势函数 $q$ 的反问题，本文提出了一种结合正则化技术的物理信息神经网络方法（Tikhonov-PINNs）。在传统 PINNs 的损失函数中引入了 Tikhonov 正则化项（$H^2$ 范数惩罚），以解决反问题的不适定性。证明了重构解关于噪声水平 $\delta$ 的收敛率（例如 $O(\delta \log(1/\delta))$），并给出了在正性条件下的具体收敛阶。在实验上通过与传统有限元方法（FEM）对比，证明了 Tikhonov-PINNs 在处理不连续势函数和高噪声数据时具有更高的精度和鲁棒性，并成功应用于 5 维反问题 。

% 第五章则利用类似技术提出了一套求解双曲型方程的算子学习算法，并给出了相应的收敛性分析


\section{预备知识}

本文主要考察如下形式的椭圆型方程（Elliptic Equations）：
\begin{equation}\label{eq:general_elliptic}
  -\Delta u(x) + f(u, x) = g(x), \quad x \in \Omega,
\end{equation}
其中 $\Omega$ 为有界开区域。当非线性项 $f(u, x)$ 关于 $u$ 呈线性依赖，即 $f(u, x) = V(x)u(x)$ 时，方程 \eqref{eq:general_elliptic} 退化为线性椭圆方程，其中 $V(x)$ 被称为势函数（Potential Function）。反之，该方程被称为非线性椭圆方程。

特征值问题（Eigenvalue Problem）通常主要关注如下齐次方程形式：
\begin{equation}\label{eq:eigen1}
  -\Delta u(x) + V(x)u(x) = \lambda u(x), \quad x \in \Omega,
\end{equation}
或者简写为 $\mathcal{L}u = \lambda u$。根据谱理论，仅当 $\lambda$ 取一系列离散实数值时，上述齐次方程存在非平凡解（Non-trivial Solution）。这些特定的 $\lambda$ 值被称为特征值，对应的解 $u(x)$ 构成的特征函数族形成了 $L^2(\Omega)$ 空间的一组正交基。

此外，本文亦关注如下形式的反问题：
\begin{equation}\label{eq:inverse1}
  -\Delta u(x) + V(x) u(x) = g(x), \quad x \in \Omega.
\end{equation}
此类问题通常定义为：在已知源项 $g(x)$ 以及解 $u(x)$ 在区域边界或内部某些观测点数据的条件下，反演重构未知势函数 $V(x)$ 的过程。这一过程被称为反势能问题（Inverse Potential Problem）。

\section{符号与定义}

\subsection{神经网络架构与数学定义}
本节将建立全连接深度神经网络的数学框架并引入相关记号。

\begin{definition}[全连接神经网络]
我们将全连接深度神经网络表示为参数化函数 $v_{\phi}: \mathbb{R}^{d} \to \mathbb{R}$。该网络通过层级递归的方式定义如下：
\begin{equation}\label{eq:nn:fcnetwork}
\begin{aligned}
v_{0}(x) &= x, \\
v_{\ell}(x) &= \sigma_{\ell}\left(A_{\ell} v_{\ell-1} + b_{\ell}\right), \quad \ell=1, 2, \cdots, L-1, \\
v_{\phi}(x) &= v_{L}(x) = A_{L} v_{L-1} + b_{L},
\end{aligned}
\end{equation}
其中，$L$ 表示网络的层数（即深度）。对于第 $\ell$ 层，$\boldsymbol{A}_{\ell} \in \mathbb{R}^{N_{\ell} \times N_{\ell-1}}$ 为权重矩阵，$b_{\ell} \in \mathbb{R}^{N_{\ell}}$ 为偏置向量，$N_{\ell}$ 表示该层的神经元个数（其中 $N_0 = d$ 为输入维数， $N_L=1$ 为输出维数）。

函数 $\sigma_{\ell}$ 为按分量的激活函数。在本文中主要用到的激活函数包括：
\begin{itemize}
  \item 线性整流单元（Rectified Linear Unit, ReLU），定义为 $\sigma(x) = \mathrm{ReLU}(x) = \max\{0, x\}$；
  \item $k$次幂次整流单元（Rectified Power Unit, RePU）定义为 $\sigma(x) = \mathrm{ReLU}^{k}(x) = \max\{0, x\}^k$，其中 $k \in \mathbb{N}^{+}$。
  \item Tanh函数：定义为 $\sigma(x) = \tanh(x) = \frac{e^{x} - e^{-x}}{e^{x} + e^{-x}}$。
\end{itemize}
\end{definition}

\begin{definition}[神经网络的参数]
网络的深度（Depth）$\mathcal{D}$ 和宽度（Width）$\mathcal{W}$ 定义为：
\begin{equation*}
\mathcal{D} = L, \quad \mathcal{W} = \max_{\ell=1, \dots, L} N_{\ell}.
\end{equation*}

此外，我们用 $\mathcal{S}$ 表示非零权重的总数。 $\mathcal{B}$ 表示非零权重绝对值的上界，即
\begin{align*}
S&=\sum_{\ell=1}^{L}(\|A_{\ell}\|_{0}+\|b_{\ell}\|_{0}),\\
B&=\max_{0\leq\ell\leq L}\max\{\max_{i,j}|A_{\ell,ij}|,\max_{i}|b_{\ell,i}|\}. 
\end{align*}

网络中所有可训练自由参数的集合记为 $\phi = \{A_{\ell}, b_{\ell}\}_{\ell=1}^{L}$。
\end{definition}

为了刻画逼近理论所需的函数空间，我们进一步定义一类满足特定复杂度约束的神经网络函数集合 $\mathcal{N}(\mathcal{D}, \mathcal{W}, \{k_{i}\}, m)$：

\begin{definition}[神经网络函数类 $\mathcal{N}$]
令 $\mathcal{N}(\mathcal{D}, \mathcal{W}, \{k_{i}\}, m)$ 表示满足以下条件的神经网络函数集合：
\begin{enumerate}
    \item {\heiti 规模约束}：网络实际深度 $D$ 和宽度 $W$ 满足 $D < \mathcal{D}$ 且 $W < \mathcal{W}$（其中 $D, W \in \mathbb{N}^+$）；
    \item {\heiti 激活函数}：所有激活函数均属于族类 $\{\mathrm{ReLU}^{k} \mid k \in \{k_{i}\}\}$；
    \item {\heiti 复杂度下界}：网络的结构参数满足如下不等式约束：
    \begin{equation}
        (D - 2 - \log_{2} W) W \geq m - 1.
    \end{equation}
\end{enumerate}
在上下文明确且不引起混淆的情况下，我们将该函数类简记为 $\mathcal{N}$。
\end{definition}

尽管全连接神经网络具有普适逼近性质，但在处理深层结构时，往往面临梯度消失或梯度爆炸的问题，导致训练困难。为了克服这一挑战，本文在实际计算中采用了残差神经网络（Residual Neural Network, ResNet）。

与方程\eqref{eq:nn:fcnetwork}定义的标准前馈结构不同，残差网络引入了“跳跃连接”（Skip Connection），使得网络学习的是恒等映射的残差。

\begin{definition}[残差神经网络]
一个包含 $L$ 个残差块（Residual Blocks）的深度残差网络 $v_{\phi}: \mathbb{R}^{d} \to \mathbb{R}$ 定义如下：

首先，通过一个线性变换（或浅层网络）将输入 $x$ 映射到隐藏层空间：
\begin{equation}
v_{0}(x) = A_{in} x + b_{in},
\end{equation}
其中 $A_{in} \in \mathbb{R}^{m \times d}, b_{in} \in \mathbb{R}^{m}$，此处 $m$ 为残差块的宽度（Width）。

随后，数据流经 $L$ 个残差块。第 $\ell$ 个残差块（$\ell=1, \dots, L$）的更新规则为：
\begin{equation}\label{eq:resnet:block}
v_{\ell}(x) = v_{\ell-1}(x) + \mathcal{F}_{\ell}(v_{\ell-1}(x)),
\end{equation}
其中 $\mathcal{F}_{\ell}$ 为残差函数。在最简单的形式中，$\mathcal{F}_{\ell}$ 由单层线性变换与激活函数构成：
\begin{equation}
\mathcal{F}_{\ell}(z) = \sigma\left(A_{\ell} z + b_{\ell}\right),
\end{equation}
这里要求 $A_{\ell} \in \mathbb{R}^{m \times m}, b_{\ell} \in \mathbb{R}^{m}$ 以保证维度匹配，从而实现逐元素相加。

最后，通过输出层将高维特征映射回目标空间：
\begin{equation}
v_{\phi}(x) = A_{out} v_{L}(x) + b_{out},
\end{equation}
其中 $A_{out} \in \mathbb{R}^{1 \times m}, b_{out} \in \mathbb{R}$。
\end{definition}

方程\eqref{eq:resnet:block}中的项 $v_{\ell-1}$ 被称为跳跃连接。这种结构允许梯度在反向传播过程中无损地流过恒等路径，从而显著增强了深层网络的训练稳定性与表达能力。

\par 接下来，我们介绍一些常用范数的符号。
\begin{definition}[范数与内积]
设 $\Omega \subset \mathbb{R}^d$ 为有界区域。对于定义在 $\Omega$ 上的函数 $u_\in L^2(\Omega)$，其 $L^2$-范数定义为：
\begin{equation*}
\|u\|_{L^2(\Omega)} = \left( \int_{\Omega} |u(x)|^2 \text{dx} \right)^{1/2}.
\end{equation*}
为简便起见，若无特殊说明，下文常将 $L^2$-范数简记为 $\|\cdot\|_{L^2}$ 或 $\|\cdot\|$。

对于任意正整数 $m \in \mathbb{N}^+$，Sobolev 空间 $H^m(\Omega)$ 中的 $H^m$-范数借由多重指标 $\alpha$ 定义为所有阶数不超过 $m$ 的弱导数之 $L^2$-范数的平方和的平方根：
\begin{equation*}
\|u\|_{H^m(\Omega)} = \left( \sum_{|\alpha| \le m} |D^\alpha u|_{L^2(\Omega)}^2 \right)^{1/2}.
\end{equation*}
特别地，常用的 $H^2$-范数具体展开为：
\begin{equation*}
\|u\|_{H^2(\Omega)}^2 = \|u\|_{L^2(\Omega)}^2 + \sum_{k=1}^d \|\partial_k u\|_{L^2(\Omega)}^2 + \sum_{k=1}^d \sum_{\ell=1}^d \|\partial_k \partial_\ell u\|_{L^2(\Omega)}^2.
\end{equation*}

针对分数阶 Sobolev 空间，函数 $u \in H^{1/2}(\Omega)$ 的范数由其 $L^2$-范数与 Gagliardo 半范数共同构成：
\begin{equation*}
\|u\|_{H^{1/2}(\Omega)}^2 = \|u\|_{L^2(\Omega)}^2 + |u|_{H^{1/2}(\Omega)}^2,
\end{equation*}
其中，衡量半阶平滑度的 Gagliardo 半范数定义为：
\begin{equation*}
|u|_{H^{1/2}(\Omega)}^2 = \iint_{\Omega \times \Omega} \frac{|u(x) - u(y)|^2}{|x - y|^{d+1}} \text{dxdy}.
\end{equation*}

上述空间皆为 Hilbert 空间，其范数均由相应的内积诱导。对应的内积分别记作 $\langle \cdot, \cdot \rangle_{L^2(\Omega)}$、$\langle \cdot, \cdot \rangle_{H^m(\Omega)}$ 与 $\langle \cdot, \cdot \rangle_{H^{1/2}(\Omega)}$，具体定义如下：
\begin{align*}
\langle u, v \rangle_{L^2(\Omega)} &= \int_{\Omega} u(x)v(x) \text{dx}, \\
\langle u, v \rangle_{H^m(\Omega)} &= \sum_{|\alpha| \le m} \langle D^\alpha u, D^\alpha v \rangle_{L^2(\Omega)}, \\
\langle u, v \rangle_{H^{1/2}(\Omega)} &= \langle u, v \rangle_{L^2(\Omega)} + \iint_{\Omega \times \Omega} \frac{(u(x) - u(y))(v(x) - v(y))}{|x - y|^{d+1}}  \text{dxdy}.
\end{align*}

若无混淆，在特定的上下文语境中，内积的下标通常会被省略，统一记作 $\langle \cdot, \cdot \rangle$。
\end{definition}